% ==========================================================================
% APPENDICES
%
% Appendix A: Hyperparameter Search and MSFT-Net Component-Addition Study
% Appendix B: Dataset Statistics and Record-Disjoint Partition Details
%
% Appendices are lettered A, B, C ... and their figures and tables are
% numbered A.1, A.2, B.1 and so on. The \appendix command below does this
% automatically.
% ==========================================================================

\appendix
\clearpage
\phantomsection
\addcontentsline{toc}{chapter}{Appendices}


% ==========================================================================
\chapter{Hyperparameter Tuning and Ablation Study Results}
\label{app:hyperparameter}
% ==========================================================================

This appendix documents the model-selection experiments behind the final
inter-patient configuration: the hyperparameter search space explored for
MSFT-Net and the component-addition study that justified its two auxiliary
components. The twelve-arm ablation of the hybrid model referenced in
Section~\ref{subsec:ablation} is reported in Table~\ref{tab:ablation} of
Chapter~\ref{ch:results} and is not repeated here. All results in this
appendix are under the inter-patient protocol; the best configuration is
reported in Chapter~\ref{ch:results}.


\section{Search Grid}
\label{app:search_grid}

Table~\ref{tab:grid} defines the search space from which 240 configurations
were sampled during the random search; the best value found for each
hyperparameter is listed.

\begin{table}[!t]
\centering
\caption{Hyperparameter search grid for the inter-patient protocol.}
\label{tab:grid}
\begin{tabular}{@{}lll@{}}
\toprule
\textbf{Hyperparameter} & \textbf{Values} & \textbf{Best} \\
\midrule
Learning rate       & $\{10^{-2},\; 10^{-3},\; 10^{-4}\}$
                    & $10^{-3}$ \\
Batch size          & $\{32,\; 64,\; 128\}$
                    & 64 \\
Stem filters        & $\{16,\; 24,\; 32\}$
                    & 32 \\
SE ratio            & $\{4,\; 8,\; 16\}$
                    & 8 \\
Dropout rate        & $\{0.2,\; 0.3,\; 0.4\}$
                    & 0.3 \\
FiLM conditioning   & \{on,\; off\}
                    & on \\
RR-interval features & \{on,\; off\}
                    & on \\
Weight decay ($L_2$) & $\{0,\; 10^{-5},\; 10^{-4}\}$
                    & $10^{-5}$ \\
\bottomrule
\end{tabular}
\end{table}


\section{MSFT-Net Component-Addition Study}
\label{app:ablation}

The FiLM conditioning and RR-interval features of MSFT-Net were selected
through a factorial component study. Starting from a baseline without
either component, each configuration enables one or both, and every
configuration was retrained and evaluated on the frozen inter-patient test
partition. The resulting macro-$F_1$ values are reported in
Table~\ref{tab:msft_addition}; each component contributes a substantial
gain on its own, and their combination is super-additive, reaching the 87.25\% macro-$F_1$ of the final MSFT-Net
(Tables~\ref{tab:perclass} and~\ref{tab:headtohead}). Per-class sensitivities and positive
predictivities for this final model are given with their Wilson confidence
intervals in Table~\ref{tab:ec57}.

\begin{table}[!t]
\centering
\caption{Effect of FiLM conditioning and RR-interval features in MSFT-Net (macro-$F_1$, inter-patient test partition)}
\label{tab:msft_addition}
\begin{tabular}{@{}llr@{}}
\toprule
\textbf{Arm} & \textbf{Configuration} & \textbf{Macro-$F_1$ (\%)} \\
\midrule
1 & Baseline (no FiLM, no RR) & 79.32 \\
2 & FiLM only                 & 82.15 \\
3 & RR features only          & 83.89 \\
4 & Both (full MSFT-Net)      & 87.25 \\
\bottomrule
\end{tabular}
\end{table}


% ==========================================================================
\chapter{Dataset Statistics and Partition Details}
\label{app:dataset}
% ==========================================================================

This appendix provides the complete database-level statistics for all
three PhysioNet databases used in this thesis, the exact
record-disjoint partition lists, and the class distribution per
partition.


\section{Database Overview}
\label{app:db_overview}

Table~\ref{tab:db_stats} summarises the source databases, their
sampling rates, record counts, and total duration.

\begin{table}[!t]
\centering
\caption{Source database statistics.}
\label{tab:db_stats}
\begin{tabular}{@{}lrrrc@{}}
\toprule
\textbf{Database} & \textbf{Records} & \textbf{Fs (Hz)} & \textbf{Total Beats} & \textbf{Duration} \\
\midrule
MIT-BIH Arrhythmia   & 48  & 360 & 109,494 & 30.5 h \\
MIT-BIH NSR          & 18  & 128 &  18,266 & 42.8 h \\
BIDMC CHF            & 15  & 250 &  15,330 & 36.0 h \\
\midrule
\textbf{Total}       & \textbf{81} & --- & \textbf{143,090} & \textbf{109.3 h} \\
\bottomrule
\end{tabular}
\end{table}


\section{Record-Disjoint Partition Lists}
\label{app:partitions}

The de~Chazal-style DS1/DS2 partition is applied only to the MIT-BIH
records. NSR and CHF records are split by record with no overlap
between partitions.

\subsection{MIT-BIH Arrhythmia Database}

\begin{table}[!t]
\centering
\caption{MIT-BIH record-disjoint partition (de~Chazal).}
\label{tab:mitbih_partition}
\begin{tabular}{@{}ll@{}}
\toprule
\textbf{Partition} & \textbf{Record Numbers} \\
\midrule
DS1 -- Train & 101, 106, 108, 109, 112, 114, 115, 116, 118, 119, \\
             & 122, 124, 201, 203, 205, 207, 209, 102, 208, 104, \\
             & 107 \\
DS1 -- Val   & 215, 220, 223, 230 \\
DS2 -- Test  & 100, 103, 105, 111, 113, 117, 121, 123, 200, 202, \\
             & 210, 212, 213, 214, 219, 221, 222, 228, 231, 232, \\
             & 233, 234, 217 \\
\bottomrule
\end{tabular}
\end{table}

\noindent
Records 102, 104, and 107 are paced records retained in DS1 for training
regularisation only; record 217 is the single held-out paced record in DS2
and supplies the entire $Q$-class test support of Table~\ref{tab:dataset}.
Consistent with Section~\ref{subsec:dataset}, paced beats contribute to the
$Q$ class only and are never scored as N/S/V/F.

\subsection{MIT-BIH Normal Sinus Rhythm Database}

\begin{table}[!t]
\centering
\caption{NSR record partition.}
\label{tab:nsr_partition}
\begin{tabular}{@{}ll@{}}
\toprule
\textbf{Partition} & \textbf{Record Numbers} \\
\midrule
Train  & nsr001--nsr012 (12 records) \\
Val    & nsr013--nsr015 (3 records)  \\
Test   & nsr016--nsr018 (3 records)  \\
\bottomrule
\end{tabular}
\end{table}

\subsection{BIDMC Congestive Heart Failure Database}

\begin{table}[!t]
\centering
\caption{CHF record partition.}
\label{tab:chf_partition}
\begin{tabular}{@{}ll@{}}
\toprule
\textbf{Partition} & \textbf{Record Numbers} \\
\midrule
Train  & chf001--chf009 (9 records)   \\
Val    & chf010--chf012 (3 records)   \\
Test   & chf013--chf015 (3 records)   \\
\bottomrule
\end{tabular}
\end{table}


\section{Class Distribution per Partition}
\label{app:class_dist}

The beat-level class distributions for each partition are reported in the
main body and are not repeated here: Table~\ref{tab:dataset} gives the
five-class inter-patient composition (train/validation/test), and
Table~\ref{tab:intra_data} gives the six-category intra-patient corpus from
raw extraction through capping to the SMOTE-balanced training set. The
corpus is severely imbalanced -- before quota sampling and balancing,
normal beats outnumber fusion beats by roughly 19:1 (15{,}306 vs.~802
beats) -- which motivated the three-way comparison of imbalance strategies
in Section~\ref{sec:workflow}. As the ablation of
Section~\ref{subsec:ablation} shows, training-only minority augmentation
with no global re-weighting was found to be the best strategy under the
record-disjoint protocol; class weighting and SMOTE were tested and
rejected.
